% !TEX root = ../main.tex
% !TEX program = lualatex
\documentclass[../main.tex]{subfiles}
\IfSubfilesClassLoaded{\externaldocument{\subfix{../build/main}}}{}
% =============================================================================
% File: 19_Fleet_Maintenance_Assets.tex
% =============================================================================
\begin{document}
\IfSubfilesClassLoaded{\chapter{EDO Study Guide}}{}
\section{Fleet Maintenance Assets}
\minibib

\subsection{Echeloned Maintenance Levels}
\begin{description}
  \item[\textbf{Organizational-level.}] Ship's force keeps the hull, mechanical, and electrical plants safe and combat-ready by executing preventative checks, battle damage control, and light repairs with organic tools, allowing the platform to remain as self-sufficient as possible while deployed~\autocite{OPNAVINST4700-7M,edo-4-1-1-fleet-maint-assets-2025}.
\item[\textbf{Intermediate-level.}] Intermediate Maintenance Activities (\acp{ima}) and tenders augment the crew with certified artisans (e.g., welders, riggers, calibration labs) who can disassemble components, correct corrosion, and return equipment to specification without docking the ship; the Navy Afloat Maintenance Training Strategy (\ac{namts}) builds those NECs, and so strike groups can recover faster when forward~\autocite{OPNAVINST4700-7M,edo-4-1-1-fleet-maint-assets-2025}.
\item[\textbf{Depot-level.}] Naval Shipyards (\acp{nsy}), private yards, and \acp{supship} conduct heavy industrial work such as reactor refueling, shaft and underwater body repairs, and major modernizations that exceed shipboard capability; every task is tied to the Class Maintenance Plan (\ac{cmp}) and engineered work instructions~\autocite{NAVSEAINST5450-14D,NAVSEAINST5450-145A,edo-4-1-1-fleet-maint-assets-2025}.
\end{description}
\note{Board cue: reference at least one example task in each echelon (e.g., \ac{ima} balancing pump rotors vs.\ depot replacement of a main reduction gear) to show command of maintenance depth.}

\subsection{Shore Maintenance Organizations and Roles}
\begin{description}
\item[\textbf{\acp{nsy}.}] Public yards (e.g., Norfolk, Pearl Harbor, Puget, Portsmouth) report to \ac{navsea} 04 and integrate nuclear workforces, planning codes, and lifting/handling departments that mirror the NAVSEAINST~5450.14D standard organization~\autocite{NAVSEAINST5450-14D}.
\item[\textbf{\acp{rmc}.}] Regional Maintenance Centers deliver voyage repairs, execute \ac{cmp} work, and act as the local \ac{ta} for availabilities; they synchronize contracting, production shops, and total ship readiness assessments (TSRA) under Commander, Navy Regional Maintenance Center (\ac{cnrmc})~\autocite{NAVSEAINST5450-145A,edo-4-1-1-fleet-maint-assets-2025}.
\item[\textbf{\ac{supship}.}] Supervisors of Shipbuilding administer new construction and conversion/modernization contracts as Administrative Contracting Officers (ACOs), coordinate budget execution, and provide technical oversight for nuclear refueling overhauls (e.g., CVN \ac{rcoh})~\autocite{NAVSEAINST5450-145A,edo-4-1-1-fleet-maint-assets-2025}.
\item[\textbf{Private shipyards.}] Huntington Ingalls, Bath Iron Works, General Dynamics NASSCO/Electric Boat, Fincantieri Marinette Marine, and Austal execute modernization packages, submarine maintenance, and surge depot work as part of the industrial base portfolio directed by \ac{opnav} and \ac{peo} contracts~\autocite{edo-4-1-1-fleet-maint-assets-2025}.
\end{description}
% TODO: Add figure that maps O-/I-/Depot-level maintenance tasks to NSYs, RMCs, SUPSHIPs, and private shipyards.

\subsection{Command, Funding, and Technical Authority Chains}
\begin{description}
\item[\textbf{Operational chain.}] \ac{cno} delegates readiness to \ac{usff} and \ac{pacflt}, who assign Type Commanders to own production schedules and risk acceptance for their forces~\autocite{OPNAVINST4700-7M}.
\item[\textbf{Technical chain.}] \ac{navsea} 21/04 issue modernization and maintenance standards, while warfare centers and \acp{ta} publish \ac{cmp}-derived work packages; disputes escalate through the Technical Authority chain independent of the operational chain~\autocite{OPNAVINST4700-7M,NAVSEAINST5450-145A}.
\item[\textbf{Funding chain.}] Fleet maintenance (OC\&MF) flows from \ac{opnav} N82 to \ac{bso}s and on to execution activities, while modernization is forecast via FUNDS~D/ALT/FUNDS~K and, in limited pilots, \ac{opn} funds for selected private-yard availabilities to accelerate contract awards~\autocite{edo-4-1-1-fleet-maint-assets-2025,DoDFMR-Vol3}.
\end{description}
Maintaining clarity between supported/supporting chains prevents the common failure mode where operational priorities outpace technical warrant approvals or funding authority, triggering rework and availability churn.

\subsection{Workforce Readiness and Assessments}
\begin{description}
\item[\textbf{Readiness assessments.}] TSRA events and Condition Found Reports feed the current ship's maintenance project (CSMP) so planners can refresh the \ac{cmp} and future availability work packages~\autocite{edo-4-1-1-fleet-maint-assets-2025}.
\item[\textbf{Training pipelines.}] Programs such as \ac{namts}, SurgeMain reservists, and civilian apprentice pipelines grow targeted skills (e.g., non-skid application, inside machinist) to close gaps identified in Fleet readiness reviews~\autocite{edo-4-1-1-fleet-maint-assets-2025}.
\item[\textbf{Digital enablers.}] Class maintenance databases, Navy Data Environment (\ac{nde}) instances, and tool control systems give \ac{rmc}s and \acp{nsy} shared configuration visibility that underpins quality assurance and auditability~\autocite{NAVSEAINST5450-145A}.
\end{description}
Timely feedback from these mechanisms is what allows \ac{surfmepp}/\ac{submepp} to rebalance work between the O-, I-, and depot echelons before a ship misses a CNO availability window.
\IfSubfilesClassLoaded{
  \RenewDocumentCommand{\entryname}{}{\textbf{\color{Modern} Acronym}}
  \RenewDocumentCommand{\descriptionname}{}{\textbf{\color{Modern} Definition}}
  \printnoidxglossary[
    type=\acronymtype,
    title=Acronyms,
    style=long-booktabs
]}{}
\end{document}
